\noindent The project introduces innovative contributions by combining XR technologies with scalable digital heritage workflows.

Key innovations include:
\begin{itemize}
    \item Use of high-fidelity digital twins as reusable cultural assets
    \item Integration of XR environments with real-world heritage preservation strategies
    \item Standardized workflows for outdoor heritage digitization
\end{itemize}

\vspace{1em}
\noindent Further breakthrough areas include:
\begin{itemize}
    \item \textbf{New XR Application Domain:} Shifting advanced 4D reality capture (recording motion, gesture, sound, and physiological data) from industrial and healthcare sectors directly into cultural heritage. This represents a highly innovative use-case that aligns with EU "cultural heritage digitisation" priorities.
    \item \textbf{Validation in Complex Scenarios:} Pushing the technological boundaries by validating scanning technologies on outdoor, deteriorating ruins under real-world conditions, rather than controlled lab environments.
    \item \textbf{Open Access Revenue and Science Models:} Implementing an innovative ecosystem where digital heritage + XR research acts as fuel for an open science platform, increasing accessibility for developers, startups, and museums.
\end{itemize}